\section{Contract design: Investor's problem}
\label{sec:contract}
\subsection{Investor's problem and equilibrium}
\label{sec:investor}

Building on the incentive-compatible mapping derived in Section~\ref{sec:environment}, the main comparative statics follow directly from the mechanism. Stronger flow convexity $f_2$ makes volatility more privately valuable to the manager because it increases upside scaling; therefore, keeping the manager at the desired risk target requires a more restrictive contract, which in the model is reflected in a lower incentive-compatible management fee. Conversely, a more risk-averse manager (higher $\eta$) is naturally disciplined and can be offered a higher fee for a given target without inducing excessive risk shifting.

We now characterize the investor's contract design problem, taking into account incentive compatibility and participation constraints. The investor chooses \(\sigma_d\) (hence \(\phi(\sigma_d)\)) to maximize its utility subject to IC and the manager's participation constraint:\footnote{We normalize the manager's reservation utility $\bar{U} = 0$. This normalization is without loss of generality: the manager's participation constraint is $U_M(\sigma_d, \phi(\sigma_d)) \ge 0$, which is verified to be slack at the calibrated equilibrium in Section~\ref{sec:numerical}. On the investor side, we similarly abstract from an outside option, though one can add the constraint $U_I\big(\sigma_d,\phi(\sigma_d)\big) \geq \bar{U}_I$ in the presence of competing managers.}
\begin{align}
\max_{\sigma_d>0}\quad & U_I\big(\sigma_d,\phi(\sigma_d)\big) \\
\text{s.t.}\quad & U_M\big(\sigma_d,\phi(\sigma_d)\big)\ge \bar U. \nonumber
\end{align}
Investor utility is given by
\begin{equation}
U_I(\sigma_d,\phi) \;=\; (1-\phi)\,(\rf+\Sharpe\sigma_d) - \phi \;-\; \frac{\gamma}{2}\,(1-\phi)^2\,\sigma_d^2.
\end{equation}
Differentiating the investor's objective with respect to $\sigma_d$ using the chain rule yields
\begin{align}
\frac{d U_I}{d \sigma_d} &= \Big[(1-\phi)\Sharpe - \gamma(1-\phi)^2\sigma_d\Big] + \Big[-(1+\rf+\Sharpe\sigma_d) + \gamma(1-\phi)\sigma_d^2\Big]\frac{d\phi}{d\sigma_d}. \label{eq:investor-FOC}
\end{align}

\begin{proposition}[Optimal Volatility Target]
Under the maintained assumptions that $S > 0$, $f_2 > 0$, and $\eta, \gamma > 0$, an interior equilibrium volatility target $\sigma_d^* > 0$ is characterized by the first-order condition \eqref{eq:investor-FOC} equal to zero. The condition reflects both the direct effect of volatility on expected returns and risk (the first bracketed term) and the indirect effect through the incentive-compatible fee adjustment $d\phi/d\sigma_d$ (the second bracketed term). The equilibrium fee is then $\phi^* = \phi(\sigma_d^*)$ from \eqref{eq:phi-IC}. In general, the resulting equation is nonlinear in $\sigma_d^*$ and is solved numerically; the comparative statics below are established by total differentiation of the first-order condition.
\end{proposition}

The first term in the first-order condition captures the standard trade-off between expected return and risk: increasing volatility raises expected returns when the Sharpe ratio is positive but also increases variance. The second term reflects the contract adjustment required to sustain incentive compatibility. Because higher target volatility requires adjusting fees to discipline the manager, changes in $\sigma_d$ affect investor utility not only through portfolio risk but also through the share of returns retained after fees.

This program summarizes the investor's contract design problem: choose a volatility target $\sigma_d$ that maximizes investor utility while guaranteeing that the manager is willing to participate and implement the target. The conflict of interests is immediate in the model. The investor values higher $\sigma_d$ only to the extent that the fund's skill (Sharpe ratio) compensates for additional risk, net of fees. The manager, in contrast, values higher $\sigma$ not only because it increases expected performance when $\Sharpe>0$, but also because it increases the probability and magnitude of inflows triggered by extreme positive outcomes. Incentive compatibility provides the key mechanism for alignment: for each target $\sigma_d$ there is a fee level that makes the manager choose exactly that risk, turning the investor's policy choice into an implementable contract. The equilibrium target and the associated fee can be computed numerically; the first-order condition and the derivative of the fee schedule are given in Appendix~\ref{app:investor}.

The structure of the optimal contract can be characterized through comparative statics without solving the model numerically. First, an increase in flow convexity $f_2$ increases the option value of volatility for the manager, hence the IC fee \(\phi(\sigma_d)\) must fall to keep incentives aligned (see Figure~\ref{fig:comp-f2}). For the investor, the contract becomes more ``expensive'' to adjust (large \(|d\phi/d\sigma_d|\)), pushing optimal \(\sigma_d\) down. Second, an increase in the scale of the agency problem (higher initial AUM $A_0$ or the linear term $f_1$) amplifies AUM sensitivity to performance; as shown in Figure~\ref{fig:comp-scale}, both desired volatility \(\sigma_d\) and the IC fee \(\phi(\sigma_d)\) tend to decrease. Third, managerial skill (an increase in the Sharpe Ratio of the fund) improves the direct mean gain from risk taking; this tends to raise the optimal target, but the IC adjustment offsets part of the effect (see Figure~\ref{fig:comp-sharpe}). Finally, as illustrated in Figure~\ref{fig:comp-risk}, an increase in the investor's risk aversion ($\gamma$) reduces \(\sigma_d\); while a higher manager $\eta$ relaxes the need for fee discipline and allows a higher \(\phi\) for a given \(\sigma_d\).

For implementation, the linear approximation \eqref{eq:phi-linear-approx} introduced in Section~\ref{sec:manager} provides a simple way to communicate the contract as a base fee plus a marginal charge for volatility. Over the relevant range of volatility targets (e.g., 6--15\% annualized), this affine rule closely tracks the exact IC fee (Appendix~\ref{app:investor}).
